\section{完整实验结果}

\subsection{统一 benchmark 全表}

统一 harness 已于 2026-06-14 完成，入口脚本为 \texttt{algos/ep11\_dl\_benchmark/scripts/run\_unified\_harness\_t1\_t2.py}，产物位于 \texttt{output/ep11\_unified\_harness/}。
完整列（runtime、source、cache path、temperature sanity、FRC@20/10 等）保存在 \texttt{all\_arm\_metrics.csv}；主文 T1/T2 分别使用 \texttt{t1\_metrics.csv} 与 \texttt{t2\_metrics.csv}。
所有 artifact 数字均为 EP11/common.metrics harness scale，不能与 TensorBoard \texttt{eval\_real/*} artifact 混表；
\texttt{tb\_vs\_harness\_scale\_check.csv} 给出三条对照（如 \eng{Stats+HP-FC} / v9b@11K: TB 0.3385 vs harness 1.7662）。

\subsubsection{T1 主表精选列}

\begin{table}[ht]
  \centering
  \caption{T1 主表精选列（统一 harness）}
  \label{tab:d-t1}
  \scriptsize
  \resizebox{\linewidth}{!}{%
  \begin{tabular}{lrrrrrrrrrr}
    \toprule
    方法 & step & split\(\downarrow\) & artifact\(\downarrow\) & corr\(\uparrow\) & FRC16\(\uparrow\) & FRC14\(\uparrow\) & FRC12\(\uparrow\) & lattice\(\downarrow\) & sharp\(\uparrow\) & FWHM\,\(\mu\)m\(\downarrow\) \\
    \midrule
    bicubic                          & ---  & 0.047 & 1.388 & 1.000 & 0.951 & 0.945 & 0.941 & 0.001 & 0.344 & 60 \\
    drizzle                          & ---  & 0.023 & 1.138 & 0.771 & 0.439 & 0.564 & 0.765 & 0.001 & 0.481 & 45 \\
    MAP-TV (\(5\times\))             & ---  & 0.024 & 2.333 & 0.438 & 0.965 & 0.955 & 0.947 & 0.002 & 0.352 & 42 \\
    TGV                              & ---  & 0.032 & 0.695 & 0.741 & 0.975 & 0.969 & 0.955 & 0.011 & 0.999 & 40 \\
    HotLoss                          & 8K   & 0.051 & 1.789 & 0.774 & 0.752 & 0.761 & 0.217 & 0.001 & 0.656 & 40 \\
    Stats                            & 15K  & 0.073 & 1.943 & 0.758 & 0.711 & 0.660 & 0.429 & 0.001 & 0.864 & 40 \\
    Stats+HP-FC                      & 11K  & 0.054 & 1.766 & 0.777 & 0.744 & 0.741 & 0.197 & 0.001 & 0.698 & 40 \\
    Stats+Full-FC                    & 7K   & 0.054 & 1.726 & 0.771 & 0.744 & 0.708 & 0.150 & 0.001 & 0.697 & 40 \\
    Hybrid                           & 10K  & 0.054 & 1.762 & 0.719 & 0.945 & 0.870 & 0.190 & 0.001 & 0.683 & 40 \\
    Hybrid+Native-FC                 & 5K   & 0.064 & 1.669 & 0.718 & 0.891 & 0.769 & 0.195 & 0.001 & 0.766 & 40 \\
    Hybrid+ResObs (\(\lambda=1.2\))  & 15K  & 0.041 & 2.726 & 0.711 & 0.986 & 0.984 & 0.979 & 0.013 & 0.968 & 40 \\
    \bottomrule
  \end{tabular}
  }
\end{table}

边界说明：\eng{MAP-TV} 是预计算 \(5\times\) anchor；\eng{TGV} 的 split/FRC 列已用独立 \eng{TGV} half-set 重建回填
（\texttt{output/ep11\_unified\_harness/tgv\_split\_frc.json}，seed 42；full-run self-check 相对 L2 \(=\) 0.0），不再是 \eng{EP16} drizzle proxy。
该 \eng{FRC} 仍是 split-consistency proxy，不是物理分辨率证据。
\eng{Hybrid+ResObs} / \eng{Hybrid} / \eng{Hybrid+Native-FC}（run ids: V10/V9A/V9C）的缓存均值约 23\,\(^\circ\)C，说明 hybrid / ResObs 推理没有漏加 drizzle base。

\subsubsection{T2 input × anchor 精选列}

\begin{table}[ht]
  \centering
  \caption{T2 输入 \(\times\) 锚定精选列}
  \label{tab:d-t2}
  \scriptsize
  \resizebox{\linewidth}{!}{%
  \begin{tabular}{llrrrrrrr}
    \toprule
    方法 & input & anchor/参数化 & step & split\(\downarrow\) & artifact\(\downarrow\) & corr\(\uparrow\) & FRC16\(\uparrow\) & FRC14\(\uparrow\) \\
    \midrule
    Stats              & 1x stats          & none                 & 15K & 0.073 & 1.943 & 0.758 & 0.711 & 0.660 \\
    Stats+PixelShuffle & 1x stats          & PixelShuffle control & 5K  & 0.050 & 1.782 & 0.739 & 0.734 & 0.692 \\
    Stats+HP-FC        & 1x stats          & HP-FC 0.1            & 11K & 0.054 & 1.766 & 0.777 & 0.744 & 0.741 \\
    Stats+Full-FC      & 1x stats          & Full-FC 0.1          & 7K  & 0.054 & 1.726 & 0.771 & 0.744 & 0.708 \\
    Hybrid             & hybrid drizzle2x  & none                 & 10K & 0.054 & 1.762 & 0.719 & 0.945 & 0.870 \\
    Hybrid+Native-FC   & hybrid drizzle2x  & Native-FC 0.1        & 5K  & 0.064 & 1.669 & 0.718 & 0.891 & 0.769 \\
    Hybrid+ResObs      & hybrid drizzle2x  & ResObs \(\lambda=1.2\) & 15K & 0.041 & 2.726 & 0.711 & 0.986 & 0.984 \\
    \bottomrule
  \end{tabular}
  }
\end{table}

\noindent（FRC12 与 sharp 列见 \texttt{t2\_metrics.csv}；Hybrid+ResObs FRC12 \(=\) 0.979、sharp \(=\) 0.968，其余变体 sharp 见 T1。）

读数：输入端引入 drizzle2x 证据能改变 fine-window 结构权衡，但不是自动保真胜利；loss-side anchor 在 \(1\times\)/full/band/legal-hybrid 变体中均未压平真实数据漂移。
\eng{Hybrid+ResObs}（V10）把 sharp/FRC 推高，同时 artifact 与 corr 变差；它是可调 trade-off，不是可认证赢家。

\subsection{FRC 完整档案}

本节归档 \eng{FRC} 分析的全部数值结果，方法学细节见 A.4。

\subsubsection{主曲线与 cutoff}

\(1/7\) cutoff 为 17.03\,\(\mu\)m（3-seed 均值曲线），逐 seed 分别为 16.17/16.17/17.03\,\(\mu\)m（std \(=\) 0.50\,\(\mu\)m），half-bit 判据给出相同值。
控制组 cutoff 为：bicubic 正控 13.58\,\(\mu\)m（未通过预期）、漂移控制 26.20\,\(\mu\)m；zero-coverage 均值 27.2\%、最大 36.2\%。
频带全表与控制组详细分析见 A.4.3 和 A.4.4。

\subsubsection{MAP-TV 前后 split-half 对照}

以 \eng{MAP-TV}（\(\sigma = 0.2\), \(\lambda = 10^{-3}\)）对 bare drizzle 做去卷积后，split-half \eng{FRC} 在 20/16/14/12/10\,\(\mu\)m 周期段分别从 0.319/0.088/0.053/0.575/0.893 提升至 0.976/0.965/0.955/0.947/0.934。
这反映的是 split-half 一致性的改善，而非光学分辨率本身的提升——bare cutoff 仍为 17.03\,\(\mu\)m，与 M2 一致。

\subsubsection{Zigzag 轮廓指标}

\eng{MAP-TV} 基准在 zigzag 剖面上的结构收益有限：中位 FWHM 从 114\,\(\mu\)m 缩窄至 100\,\(\mu\)m，dip 从 0.929 微升至 0.934，3/3 剖面保持分离但改善不均（2 条变宽、1 条显著变窄）。
论文措辞限定为 "limited contour enhancement"。

\subsection{漂移演化档案}

\subsubsection{全局漂移指标}

下表记录了各变体在训练期间真实数据 eval（TB-scale，248 帧 contour\_refined）上的 artifact（\(\downarrow\)）与 raw\_control\_corr（\(\uparrow\)）轨迹（每格为 artifact / corr）。

\begin{table}[ht]
  \centering
  \caption{各变体真实数据 eval 漂移轨迹（artifact / corr）}
  \label{tab:drift-archive}
  \small
  \resizebox{\linewidth}{!}{%
  \begin{tabular}{lccccc}
    \toprule
    step & Stats & Hybrid & Stats+HP-FC & Stats+Full-FC & Hybrid+Native-FC \\
    \midrule
    10K & 0.390 / 0.756 & 0.446 / 0.719 & 0.369 / 0.758 & 0.379 / 0.758 & 0.516 / 0.714 \\
    20K & 0.476 / 0.729 & 0.514 / 0.702 & 0.486 / 0.735 & 0.575 / 0.642 & 0.689 / 0.655 \\
    30K & 0.602 / 0.703 & 0.660 / 0.663 & 0.611 / 0.709 & 0.615 / 0.694 & 0.686 / 0.675 \\
    40K & 0.627 / 0.698 & 0.656 / 0.665 & 0.640 / 0.697 & 0.672 / 0.681 & 0.688 / 0.672 \\
    60K & 0.643 / 0.689 & 0.646 / 0.669 & 0.655 / 0.688 & 0.677 / 0.677 & 0.695 / 0.669 \\
    \bottomrule
  \end{tabular}
  }
\end{table}

从轨迹形状看，\eng{Hybrid}（V9A）是唯一在 30K\(\rightarrow\)60K 段漂移趋平的变体（变化量 \(-0.014\)/\(+0.007\)，其余单调恶化）。
\eng{Stats+Full-FC}（V9D）比 \eng{Stats+HP-FC}（V9B）更差且 1K--28K 段震荡剧烈。
\eng{Hybrid+Native-FC}（V9C）在 30K 后没有恢复，60K 端点为 0.695/0.669，说明 hybrid 输入下的合法 \(1\times\) anchor 也不能压平漂移。
需要强调的是，hybrid 变体（\eng{Hybrid} / \eng{Hybrid+Native-FC}）与 \(1\times\) 变体的 proxy 不可跨列横比（见 A.3.2），此处只看各变体轨迹形状。

\subsubsection{Fine-window 训练时间轴}

为更精细地诊断训练动态，在 \(2\times\) 网格全幅 \(960 \times 1280\) 的中心细线窗口（rows 384:518, cols 478:674——中心两条细 zigzag 及周边粗线）上，跟踪四个指标：
\(\text{hp\_corr\_input}\)（highpass 域对 drizzle 输入通道的 Pearson，衡量保真性）、\(\text{hp\_corr\_tgv}\)（同口径对 \eng{TGV} 的 Pearson，交叉验证）、
\(\text{sharp\_p95}\)（温度图 P95 梯度幅值，衡量锐度）、\(\text{lattice\_score}\)（highpass 域 \(|f| > 0.35\) cyc/px 频段能量占比，捕捉格纹伪影）。

两个参照点为：drizzle 输入通道 \((1.000, 0.503)\)——观测域上限（模糊但零幻觉）；\eng{TGV} \((0.960, 0.959)\)——经典前沿工作点。

\eng{Hybrid}（V9A）的训练时间轴揭示了一个关键现象——30K 保真悬崖。\eng{Hybrid} 10K 的 \(\text{hp\_corr\_input}\) 为 0.970，20K 达峰 0.974，
此后断崖下跌：30K 降至 0.908 并在 40K--60K 焊死于 \(0.906 \pm 0.001\)。
同时，锐度 \(\text{sharp\_p95}\) 在 30K 后超过 \eng{TGV}（1.147 vs 0.959），并在 40K--60K 继续攀升至 1.21--1.25——锐度超过 \eng{TGV} 的区间恰好与去相关重合，指示过冲为幻觉驱动。
对照 \eng{Stats}（v8.1a）60K 同窗口为 0.926/0.936，\eng{Hybrid} 60K（0.925/0.935）与之无差别——hybrid 早期增益被训练后期完全抹平。

需注意 \eng{Hybrid}（V9A）的 35K 中断（batch size 从 128 切换至 64）与悬崖时间重合，混杂尚未排除（见 C.3）。

Fine-window 是局部诊断窗口，依赖 \eng{TGV} 参照，独立性低于 FRC/proxy 对——仅用于归因与选型，不替代全局评估协议。

\subsection{负结果档案}

本节完整记录四项负结果，为主文~\S5.7 提供数据背书。每项按「现象 \(\rightarrow\) 数字 \(\rightarrow\) 结论边界」结构组织。

\subsubsection{PixelShuffle HR 头}

\eng{Stats+PixelShuffle}（v8.1b）使用 pixelshuffle \eng{HR} 头替代默认的 bilinear 头（其余与 \eng{Stats} / v8.1a 相同），
出现了中等边框间条纹状亮色伪影，锯齿未改善，中心细线模糊程度与 bilinear 头相同。
TB-scale 数据显示 artifact 全程高于 \eng{Stats}（0.413\(\rightarrow\)0.709 vs 0.390\(\rightarrow\)0.643），corr 全程更低（0.747\(\rightarrow\)0.667 vs 0.758\(\rightarrow\)0.689）。

结论边界：head 归因失败——中心细线瓶颈不在解码头。该结果与 \eng{Stats}（v8.1a）的 loss-cooldown 对照共同构成主文~\S5.3 的两个对照归因：细线模糊对 loss 温度和 head 均不敏感，指向输入信息瓶颈。

\subsubsection{4x 网络}

\eng{EP12} \(4\times\) 网络在所有指标上均劣于 \eng{EP07} \(2\times\) 上采样方案：raw-control highpass Pearson 0.223 vs 0.389，artifact 0.535 vs 0.472。

该实验结果与 A.1.3 的理论预测一致：\(4\times\) Nyquist \eng{MTF} \(\le 0.042\)（\(\sigma = 0.2\)），\(\sigma \ge 0.35\) 时接近零；相位 bin 在 \(3\times\)/\(4\times\) 网格出现 collapse（见 B.2.3）。
\(4\times\) 无真实增益且理论上不应有——这是「先第一性原理后实验」工作流的展示案例。

\subsubsection{Forward 锚定（Stats+HP-FC / Stats+Full-FC / Hybrid+Native-FC）}

这些锚定变体的漂移曲线与无锚 \eng{Stats}（v8.1a）几乎重合。\eng{Stats+HP-FC}（V9B）的 forward loss 自 10K 步起贴底于 0.004--0.009，同期 artifact 持续上爬。
\eng{Stats+HP-FC} 40K\(\rightarrow\)60K 漂移 \(+0.0145\)/\(-0.0082\) vs \eng{Stats} \(+0.016\)/\(-0.009\)，几乎完全重合。
\eng{Stats+Full-FC}（V9D）60K（0.677/0.677）劣于 \eng{Stats+HP-FC}（0.655/0.688），且 1K--28K 段剧烈震荡——复现了全频低通梯度冲突的已知问题。

\eng{Hybrid+Native-FC}（V9C）是最后一条反对意见检验：在 hybrid 输入含 \(2\times\) drizzle 证据的条件下，forward anchor 不再偷用上采样 mean，而是通过数据管线携带合法的 \(1\times\) aligned-mean patch。
结果仍为 0.516/0.714 @10K \(\rightarrow\) 0.695/0.669 @60K，未压平后期漂移。

结论边界：无论频带选择和输入证据路径如何，\(1\times\) 观测域锚定对真实漂移不可见——机制即 A.2.3 Proposition 1。
\eng{Stats+Full-FC}（V9D）关闭了「band 太窄」的反对意见；\eng{Hybrid+Native-FC}（V9C）关闭了「hybrid 第 0 通道不是合法 \(1\times\) 观测」的反对意见。
loss-side anchoring route 在 HP-FC / Full-FC / Native-FC 三种变体下正式关闭。

\subsubsection{AVI 作 SR 输入}

8-bit 渲染视频、约 67\% 重复帧、无温度矩阵（见 B.1.3）——直接排除。\eng{AVI} 仅用于方向一致性检验。
该负结果的意义在于：「能拿到的数据」不等于「能进重建的数据」，审计应先于算法。

\subsection{帧预算与鲁棒性}

本节报告 \eng{EP16} 的三组经典方法实验：E1 帧数预算、E2 shift 扰动、E3 对齐源消融。共 37 个 unique run 全部成功。
注意这是推理期稳定性研究，不属于统一 harness。该轨迹图原为主文 F7，现移至此处仅作 supp 参照。

\begin{figure}[ht]
  \centering
  \includegraphics[width=\linewidth]{../../output/paper_figures/fig07_budget_robustness.png}
  \caption{帧数预算（左）与 shift 扰动鲁棒性（右）。\eng{Drizzle} 的 raw-control corr 多数增益在 \(N = 62\) 即出现，
  FRC/split-half proxy 随更多相位覆盖继续改善；\eng{TGV} 全程 artifact 低于 drizzle 但 corr 非单调。
  横轴对应 \(N \in \{31,62,124,248\}\) 与 \(\sigma \in \{0,0.05,0.1,0.2\}\)\,px，详见下表。}
  \label{fig:d-budget}
\end{figure}

\subsubsection{帧数预算}

帧数预算实验使用 phase-stratified 子集（\(N \in \{31, 62, 124, 248\}\)，seeds 101/202/303），覆盖 drizzle 和 \eng{TGV} 两种经典方法。

\begin{table}[ht]
  \centering
  \caption{帧数预算扫描（drizzle 与 TGV）}
  \label{tab:budget}
  \small
  \resizebox{\linewidth}{!}{%
  \begin{tabular}{lrrrrrr}
    \toprule
    方法 & \(N\) & corr & split-half NRMSE & artifact & FRC@16\,\(\mu\)m & zigzag FWHM (\(\mu\)m) \\
    \midrule
    drizzle & 31  & 0.747\(\pm\)0.032 & 0.0715\(\pm\)0.0012 & 1.649\(\pm\)0.041 & 0.109\(\pm\)0.065 & 56.7\(\pm\)20.2 \\
    drizzle & 62  & 0.772\(\pm\)0.016 & 0.0514\(\pm\)0.0034 & 1.536\(\pm\)0.044 & 0.248\(\pm\)0.070 & 56.7\(\pm\)20.2 \\
    drizzle & 124 & 0.770\(\pm\)0.010 & 0.0361\(\pm\)0.0010 & 1.341\(\pm\)0.038 & 0.332\(\pm\)0.018 & 45.0 \\
    drizzle & 248 & 0.771 & 0.0306 & 1.145 & 0.479 & 45.0 \\
    TGV & 31  & 0.728\(\pm\)0.053 & --- & 0.946\(\pm\)0.003 & --- & 42.5\(\pm\)3.5 \\
    TGV & 62  & 0.754\(\pm\)0.009 & --- & 0.851\(\pm\)0.051 & --- & 40.0 \\
    TGV & 124 & 0.735\(\pm\)0.013 & --- & 0.747\(\pm\)0.011 & --- & 40.0 \\
    TGV & 248 & 0.741 & --- & 0.708 & --- & 40.0 \\
    \bottomrule
  \end{tabular}
  }
\end{table}

\eng{TGV} 的 split-half/FRC 列使用同子集 drizzle proxy（预算考虑），以 "---" 标注。

Drizzle corr 的增益大半在 \(N = 62\) 前到位（0.747\(\rightarrow\)0.772），之后趋平；split-half、artifact 和 \eng{FRC} 随 \(N\) 单调改善，受益于相位覆盖的累积。
\eng{TGV} artifact 全程低于 drizzle（0.946\(\rightarrow\)0.708），但 corr 呈非单调变化（0.728\(\rightarrow\)0.754\(\rightarrow\)0.735\(\rightarrow\)0.741）。

\subsubsection{Shift 扰动}

扰动实验在 \(\sigma \in \{0, 0.05, 0.1, 0.2\}\)\,px（seeds 401--403）下测试对齐误差的影响。
Drizzle corr 从 0.771 到 0.770 几乎不变，而 artifact 从 1.145 升至 1.434、FRC@16 从 0.479 降至 0.340。\eng{TGV} corr 同样稳定（0.741\(\rightarrow\)0.744）。

结论：鲁棒性是 metric-specific 的——raw-control corr 对 \(\le 0.2\)\,px 扰动几乎不敏感，而 coverage/FRC 类指标敏感。这是压力测试口径，非真实对齐误差估计。

\subsubsection{对齐源消融}

将 command\_prior 与 contour\_refined 两种对齐源的端到端效果作比：drizzle corr 从 0.662 升至 0.771（\(+0.109\)），FRC@16 从 0.0166 升至 0.479；\eng{TGV} corr 从 0.642 升至 0.741（\(+0.099\)）。

数据驱动对齐对两个经典方法都带来约 \(+0.10\) corr 的端到端增益，FRC@16 从几乎为零到 0.479——这是 B.2 对齐链投入的最终回报证据。

\subsection{视觉 gate panel}

已有 \eng{EP11} 四个变体（\eng{HotLoss}、\eng{Stats}、\eng{Stats+PixelShuffle}、\eng{Stats+HP-FC}）的 checkpoint 选择 panel：每个变体 canonical \(+\) 60K 对照的温度域三联图，标注 proxy 值。
\eng{Stats+PixelShuffle} 行保留条纹伪影证据。主文 F5 已补 \eng{Hybrid} late endpoint 60K 与 \eng{Hybrid+ResObs} \(\lambda=1.2\)@15K 的双域视觉 gate，
并把中心 \eng{ROI} 与 held-out ROI2 合成 4 行图：\texttt{output/paper\_figures/fig05\_combined\_visual.\{png,pdf\}}。
\eng{Hybrid+Native-FC} 的 step 序列 panel 可作为 supp 选图补充，不是主文硬门槛。

\eng{Hybrid} 输入变体（V9A）的单独漂移轨迹（原主文 companion，因 proxy 跨输入模式不可横比移至此处）见图~\ref{fig:d-v9a}。

\begin{figure}[ht]
  \centering
  \includegraphics[width=\linewidth]{../../output/paper_figures/fig03s_v9a_trajectory.png}
  \caption{\eng{Hybrid} 输入漂移轨迹（companion；run id: V9A）。artifact / corr vs 训练步数，纵轴尺度仅在 \eng{Hybrid} 自身轨迹内可比，不能与 \(1\times\) 输入变体横向比较。
  \eng{Hybrid} 非单调、晚期压平但平台仍低于早期 checkpoint——输入端证据注入改变漂移形态、不消除先验侵蚀。}
  \label{fig:d-v9a}
\end{figure}

\subsection{零训练融合 baseline}

本节建立一个无需额外训练的简单 baseline：对经典重建（锚）与学习方法输出做事后线性融合，测试学习方法的增益是否超出后处理可得的范围。

\subsubsection{方法}

融合公式为 \(\text{fused}(\lambda) = (1-\lambda) \cdot \text{anchor} + \lambda \cdot \hat{x}_{\text{unet}}\)，\(\lambda \in \{0, 0.1, \ldots, 1.0\}\)。
锚选择 drizzle \(2\times\) mean（248 帧 contour\_refined）或 \eng{TGV}。UNet 预测使用 \eng{Hybrid}（V9A）20K（最保真）和 \eng{Hybrid}（V9A）60K（最锐）。评估指标为 D.3.2 的 fine-window 四指标。

E3 追加了第二个固定验证窗，入口为 \texttt{algos/ep07\_unet\_sr/scripts/v9\_review/run\_fusion\_window2.py}，运行时显式设置 \texttt{CUDA\_VISIBLE\_DEVICES=""} 且只读取缓存全幅预测。
原 D.7 fine-window 为 \(2\times\) 网格 rows 384:518, cols 478:674；第二窗保持同尺寸与同列，向下平移一个窗口高度并留 24\,px 间隔，即 rows 542:676, cols 478:674。
该规则在重选 \(\lambda\) 前固定，用于检查原窗局部选择是否可迁移到另一个非重叠 \eng{ROI}。

\subsubsection{结果}

原 fine-window 上，\eng{TGV} 锚 \(\times\) \eng{Hybrid}-60K（V9A）的低 \(\lambda\) 组合形成一段局部 proxy-frontier。
\(\lambda = 0.2\) 时四指标为 hp\_corr\_input \(=\) 0.963（\eng{TGV} 参照 0.960）、hp\_corr\_tgv \(=\) 0.995、sharp\_p95 \(=\) 0.968（\eng{TGV} 参照 0.959）、lattice \(=\) 0.0108（\eng{TGV} 参照 0.0169）。
这说明在该局部窗口内，少量 \eng{Hybrid}-60K 高频内容可以改善 fine-window proxy 读数；但该读数依赖 \eng{TGV}/highpass 参照与同窗 \(\lambda\) 选择，不是 GT 保真或物理分辨率证据。

第二窗结果位于 \texttt{output/ep07\_v9\_review/fusion\_window2\_metrics.csv} 与 \texttt{fusion\_window2\_summary.md}。
第二窗局部参照为 drizzle \(=\) (hp\_corr\_input 1.000, sharp\_p95 0.427, lattice 0.0036)，\eng{TGV} \(=\) (0.962, 0.491, 0.0160)。
按同一选择规则（\eng{TGV} 锚 \(\times\) \eng{Hybrid}-60K、\(\lambda > 0\)，要求 hp\_corr\_input 与 sharp\_p95 不低于本窗 \eng{TGV}、lattice 不高于本窗 \eng{TGV}，并以 hp\_corr\_input 优先排序），
重选得到 \(\lambda = 0.1\)：hp\_corr\_input \(=\) 0.964、hp\_corr\_tgv \(=\) 0.998、sharp\_p95 \(=\) 0.508、lattice \(=\) 0.0126。
\(\lambda = 0.2\) 也通过同一 gate（0.964 / 0.994 / 0.525 / 0.0100），但因 hp\_corr\_input 略低而排第二。

相比之下，drizzle 锚 \(\times\) \eng{Hybrid} 的任意 \(\lambda\) 都没有形成同类前沿点：drizzle 参照本身保真极高，但锐度始终不足。

\eng{Hybrid+ResObs}（V10）高-\(\lambda\) 工作点的相对位置跨窗保持一致：原窗 \(\lambda = 1.2\) @15K 为 hp\_corr\_input \(=\) 0.922、sharp\_p95 \(=\) 0.987、lattice \(=\) 0.0141；
第二窗为 hp\_corr\_input \(=\) 0.920、sharp\_p95 \(=\) 0.501、lattice \(=\) 0.0130。
两窗中 \eng{Hybrid+ResObs} 均低于对应 \eng{TGV} 的 hp\_corr\_input（0.960/0.962），同时 lattice 低于 \eng{TGV}、sharp\_p95 接近或略高于 \eng{TGV}。
这支持“\eng{Hybrid+ResObs} 是低 grain / 较锐但保真不足的折中点”这一 proxy 位置判断，而不是可认证胜出。

\subsubsection{结论与边界}

第一，零训练线性融合是一个必须报告的强 sanity baseline：在两个非重叠局部窗口中，\eng{TGV} 锚 \(\times\) \eng{Hybrid}-60K 的低 \(\lambda\) 点都能进入相对 \eng{TGV} 的 proxy-frontier 区域。
第二，最优 \(\lambda\) 不完全跨窗稳定（0.2 \(\rightarrow\) 0.1），因此不能把单窗最优点写成稳定算法结论；更稳妥的表述是“低 \(\lambda\) TGV-anchored 融合在局部 proxy 上构成对学习方法的挑战”。

\eng{Hybrid+ResObs} 高-\(\lambda\) sweep 的最佳工作点仍按 \(\lambda = 1.2\) @15K 报告：原窗 hp\_corr\_input \(=\) 0.922、sharp\_p95 \(=\) 0.987、lattice \(=\) 0.0141；第二窗 hp\_corr\_input \(=\) 0.920、sharp\_p95 \(=\) 0.501、lattice \(=\) 0.0130。
它满足“低 grain、锐度不低于 \eng{TGV} 局部读数”的 proxy 判据，但保真仍低于 \eng{TGV}，因此不构成可认证胜出。

边界必须保留：fine-window 与第二窗都是局部 \eng{ROI}，指标依赖 \eng{TGV}/highpass 参照而非 GT；\(\lambda\) 是在同一真实数据窗口上重选，仍有 selection-on-test caveat；
\texttt{sharp\_p95} 不度量轮廓连续性且可受 beading/grain 影响。D.7 的结论只用于 proxy 稳定性与 baseline 压力测试，不能外推为温度计量、物理分辨率或单一方法胜出声明。
